\documentclass[10pt,a4paper]{amsart}
\usepackage[margin=19mm]{geometry}
\usepackage{amsmath,amssymb,mathtools}
\usepackage[scr=rsfs,cal=euler]{mathalfa}
\usepackage{xfrac}
\usepackage[unicode,hidelinks,bookmarks=false]{hyperref}
\newcommand{\ie}{i.e.\ }
\newcommand{\cf}{cf.\ }
\newcommand{\resp}{respectively\ }
\newcommand{\Subsection}{\subsection}
\providecommand{\nobreakdash}{\nobreak\hbox{-}}
\setlength{\emergencystretch}{2em}
\multlinegap0pt
\usepackage{bm}
\usepackage{bbm}
\usepackage{dsfont}
\usepackage{xspace}
\usepackage{cancel}
\usepackage{mathtools}
\usepackage{amsthm}
\makeatletter
\@ifundefined{theorem}{\newtheorem{theorem}{Theorem}}{}
\makeatother
\def\Ap{\operatorname{Ap}}
\def\N{\mathbb{N}}
\def\supp{\operatorname{supp}}
\title[A note on strong affine semigroups]{A note on strong affine semigroups}
\author{I. Garc\'ia-Marco, R. Tapia-Ramos,, A. Vigneron-Tenorio}
\date{Source: arXiv:2606.28420v1 (CC BY 4.0)}
\begin{document}
\maketitle

\noindent\emph{Source and licence.} A note on strong affine semigroups. Version arXiv:2606.28420v1 (CC BY 4.0), distributed under the Creative Commons Attribution 4.0 International licence, as recorded in the arXiv licence metadata for this exact version and shown on the abstract page https://arxiv.org/abs/2606.28420v1. Full rights notice: licenses/arxiv-2606.28420-CC-BY-4.0.txt.\par\smallskip

\noindent\textit{Editorial scope.} Counted unit: the Theorem whose source label is \texttt{TheoremYiYj} in the article (source0.tex lines 560--570, source label \texttt{TheoremYiYj}), delivered together with the article's own complete original proof of it (source0.tex lines 572--597, one \texttt{proof} environment of 26 lines). No mathematics is written, repaired or re-derived: statement and proof are the article's own text line for line. Required context retained uncounted: none. Every definition, display and label of those lines is retained unchanged. Omitted from this package and not redistributed: 1--559, 571--571, 598--801. Nothing inside a retained excerpt is omitted or altered: every retained body is delivered exactly as the article's lines are written, so no omission receipt of any kind accompanies this package. Citations: the retained excerpts carry no citation command at all, so no bibliography entry of the article is transcribed and no reference is redistributed. Reference adaptation: none was needed and none was made; every reference inside the delivered unit resolves inside it, so no unresolved reference or citation is produced. Printed slips kept literal: no printed word, symbol, hypothesis or number of the counted statement or of its proof was altered. Layout and class: the article's own class is \texttt{article} with packages including graphicx, amsfonts, amssymb, amsthm, amsmath, bm, ebezier, url, xcolor, algorithm2e, algpseudocode, enumitem, hyperref; the wrapper is the lane's pinned \texttt{amsart} editorial wrapper at 10pt on A4, which restates the theorem-like environments the retained text needs and adds the article's own macro definitions that the retained text needs (\texttt{\textbackslash Ap}, \texttt{\textbackslash N}, \texttt{\textbackslash supp}), with extra wrapper packages (bm, bbm, dsfont, xspace, cancel, mathtools). The delivered numbering is the unit's own. Assets: no retained span contains a figure environment, a graphics-inclusion command, a PDF-page inclusion, a table or a TikZ picture; no image file is copied into this package, the package declares no asset dependency and no image file is redistributed with it. Licence: version arXiv:2606.28420v1 of this article is distributed under the Creative Commons Attribution 4.0 International licence, as recorded in the arXiv licence metadata for this exact version and shown on the abstract page https://arxiv.org/abs/2606.28420v1. A full rights notice is delivered at licenses/arxiv-2606.28420-CC-BY-4.0.txt. Characters: every retained excerpt is pure ASCII, so the delivered engine, XeLaTeX with its default Latin Modern text font, reports no missing character.
\begin{theorem}\label{TheoremYiYj}
Let $S$ be a strong affine semigroup minimally generated by $E\sqcup A$ with $A\neq \emptyset$. Then, the following monomials are indispensable:
\begin{enumerate}
\item $y_i y_j$ for $i,j\in [r]$ with $i\neq j$.
\item $y_i^2$ for any $i\in [r]$ such that $2a_i\notin \Ap(S,E)$.
\item $y_i^3$ for any $i\in [r]$ such that $2a_i\in \Ap(S,E)$.
\item $X^\omega y_j$  for any $j\in [r]$ and $\omega \in {\rm Hilb} (T_j)$.
\item $X^\omega y_j^2$  for any $j\in [r]$ with $2a_j\in\Ap(S,E)$, and $\omega \in {\rm Hilb} (T'_j)$.
\end{enumerate}
Moreover, they are the unique indispensable monomials involving at least one variable in $\{y_1,\ldots, y_r\}$.
\end{theorem}

\begin{proof}
Let $M$ be a monomial and denote $m = S\text{-degree}(M)$. We have that $M$ is indispensable if and only if $\nabla_m$ is not connected and $\{M\}$ is a connected component. Equivalently, $M$ is indispensable if and only if the cardinality of $C_{m}$ is at least two, and for any monomial $M' \in C_m$ distinct from $M$, then $\gcd(M, M') = 1$.

Consider $i,j\in [r]$ with $i\neq j$. Since $S$ is a strong semigroup, the cardinality of $C_{a_i+a_j}$ is at least two, let $X^\alpha Y^\beta \in C_{a_i+a_j}\setminus\{y_iy_j\}$. Note that $i,j\notin \supp (\beta)$. Otherwise, $a_i$ or $a_j$ is not a minimal generator of $S$, or $S$ is such that $S\cap (-S)\neq \{0\}$. Hence, $\gcd(X^\alpha Y^\beta,y_iy_j) = 1$.
Analogously, for any $2a_i\notin\Ap(S,E)$, there also exists more than one monomial in $C_{2a_i}$, but there is no one as $y_iX^{\alpha'} Y^{\beta'}$ in $C_{2a_i}\setminus\{y_i^2\}$.
For every $y_i^3$ such that $2a_i \in \Ap (S,E)$, we know that $\sharp C_{3a_i} \ge 2$. Since $C_{2a_i}=\{y_i^2\}$, there are no monomials $y_iX^\alpha Y^\beta$ in $C_{3a_i}$ with $\alpha\neq 0$.
To prove the fourth item, we consider an integer $j\in [r]$ and an element $\omega \in {\rm Hilb} (T_j)$. Hence, $\sharp C_{\sum_{i=1}^t \omega_i e_i+a_j}\ge 2$, and $\sharp C_{\sum_{i=1}^t \omega_i a_i}=1$. Let $X^\alpha Y^\beta$ be a monomial in $C_{\sum_{i=1}^t \omega_i e_i+a_j}\setminus \{X^\omega y_j\}$. If $j \in \supp(\beta)$, then $\{X^\omega,X^\alpha \frac{Y^\beta}{y_j}\}\subseteq C_{\sum_{i=1}^t \omega_i e_i}$, but this contradicts with the cardinality of $C_{\sum_{i=1}^t \omega_i e_i}$. Furthermore, if we assume that there exists some $k\in \supp(\omega)\cap \supp(\alpha)$, then $(\omega_1,\ldots, \omega_{k-1},\omega_{k}-1,\omega_{k+1},\ldots \omega_t)$ belongs to $T_j$, but this fact implies that $\omega \notin {\rm Hilb} (T_j)$, which is not possible. Consequently, for every $X^\alpha Y^\beta \in C_{\sum_{i=1}^t \omega_i e_i+a_j}\setminus \{X^\omega y_j\}$, we have that $\gcd(X^\omega y_j,X^\alpha Y^\beta) = 1$. The last item can be proved similarly, taking $\omega \in {\rm Hilb} (T'_j)$ and assuming that $2a_j\in\Ap(S,E)$. Hence, the monomials appearing in the different items are indispensable.

Let us prove now that there are no other indispensable monomials. We consider $X^\alpha Y^\beta \in C_m$ an indispensable monomial  involving at least one variable in $\{y_1,\ldots,y_r\}$. We study different cases based on the cardinality of the support of $\beta=(\beta_1,\ldots,\beta_r)$:

\noindent {\em Case 1:} $\sharp \supp (\beta)\ge 2$. Without loss of generality, we can assume that $\beta_1\beta_2\neq 0$. Since $S$ is a strong semigroup, there exist $k\in [t]$, $\alpha'\in \N^t$, and $\beta'\in \N^r$ such that the monomial $M := x_kX^{\alpha+ \alpha'}\frac{Y^{\beta + \beta'}}{y_1y_2} \in C_m$. As $X^\alpha Y^\beta$ is indispensable, then $1 = \gcd(X^\alpha Y^\beta, M).$ Thus, $X^\alpha Y^\beta = y_1 y_2$.

\noindent {\em Case 2:} If $\sharp \supp (\beta) = 1$. Without loss of generality, we can assume that $\beta_1 \neq 0.$

{\em Case 2.1}: assuming that $2a_1 \in \Ap(E,S)$, we distinguish the following cases,

\begin{itemize}
\item If $\beta_1 \geq 3$, then  there exist $k\in [t]$, $\alpha'\in \N^t$, and $\beta'\in \N^r$ such that $3 a_1 = e_k + \sum_{i = 1}^t \alpha_i' e_i + \sum_{i = 1}^r \beta_i' a_i$. Hence, $M = x_k X^{\alpha + \alpha'} y_1^{\beta_1 - 3} Y^{\beta'}\in C_m$. Since $X^\alpha Y^{\beta}$ is indispensable, this means that $1 = \gcd(X^\alpha Y^\beta, M) = X^{\alpha}.$ Thus, $X^\alpha Y^\beta = y_1^3$.

\item  If $\beta_1 = 2$, let us prove that $\alpha \in {\rm Hilb} (T'_j)$. Assume otherwise that $\alpha \notin {\rm Hilb} (T'_j)$, this means that either: (i) $\alpha \notin T'_j$, or (ii) there exists $\alpha' \in T'_j$ such that $X^{\alpha'}$ divides $X^{\alpha}$. In (i), the set $C_{m-a_1}$ has two different elements $N_1$ and $N_2$. So, $N_1 y_1, N_2 y_1 \in C_{m}$ and $\gcd(N_iy_1,X^\alpha y_1^2) \neq 1$ for $i = 1,2$; but this contradicts that $X^{\alpha} Y^{\beta}$ is indispensable. In (ii), the set $C_{\sum_{i = 1}^t \alpha_i'e_i + 2 a_1}$ has at least two elements $N_1, N_2$. So  $X^{\alpha - \alpha'}N_1, X^{\alpha - \alpha'} N_2 \in C_{m}$ and $\gcd(X^{\alpha - \alpha'} N_i,X^\alpha y_1^2) \neq 1$ for $i = 1,2$; but this also contradicts that $X^{\alpha} Y^{\beta}$ is indispensable.

\item  If $\beta_1 = 1$, one can prove that $\alpha \in {\rm Hilb} (T_j)$ analogously.
\end{itemize}

{\em Case 2.2}: when $2a_1 \notin \Ap(E,S)$, and proceeding as in {\em Case 2.1}, one can prove that if $X^{\alpha}Y^{\beta}$ is indispensable, then $\beta_1 = 1$ and $\alpha \in {\rm Hilb} (T_j)$.
\end{proof}

\end{document}
